\hypertarget{fargate-crawler-ecs-fargate-eventbridge-scheduler}{%
\section{Fargate Crawler: ECS Fargate + EventBridge
Scheduler}}

This section covers deploying the Scrapy crawler as an ECS Fargate task
triggered by EventBridge Scheduler.

\hypertarget{architecture}{%
\subsection{Architecture}}

\emph{(Mã nguồn chi tiết đã được lược bỏ để báo cáo ngắn gọn. Vui lòng
xem trong source code dự án)}

\hypertarget{task-definition-terraform}{%
\subsection{Task Definition
(Terraform)}}

From \texttt{main.tf}:

\emph{(Mã nguồn chi tiết đã được lược bỏ để báo cáo ngắn gọn. Vui lòng
xem trong source code dự án)}

\hypertarget{eventbridge-scheduler-terraform}{%
\subsection{EventBridge Scheduler
(Terraform)}}

\emph{(Mã nguồn chi tiết đã được lược bỏ để báo cáo ngắn gọn. Vui lòng
xem trong source code dự án)}

\hypertarget{iam-permissions-for-task-role}{%
\subsection{IAM Permissions for Task
Role}}

The \texttt{ecs\_task\_role} needs permissions for: - \textbf{SQS}:
SendMessage to \texttt{news\_raw} queue - \textbf{Secrets Manager}:
GetSecretValue for DB credentials - \textbf{CloudWatch Logs}:
CreateLogStream, PutLogEvents - \textbf{ECR}: BatchGetImage,
GetDownloadUrlForLayer (via execution role)

\emph{(Mã nguồn chi tiết đã được lược bỏ để báo cáo ngắn gọn. Vui lòng
xem trong source code dự án)}

\hypertarget{sqs-queue-terraform}{%
\subsection{SQS Queue (Terraform)}}

\emph{(Mã nguồn chi tiết đã được lược bỏ để báo cáo ngắn gọn. Vui lòng
xem trong source code dự án)}

\hypertarget{building-pushing-docker-image}{%
\subsection{Building \& Pushing Docker
Image}}

\emph{(Mã nguồn chi tiết đã được lược bỏ để báo cáo ngắn gọn. Vui lòng
xem trong source code dự án)}

\hypertarget{manual-testing}{%
\subsection{Manual Testing}}

\hypertarget{run-task-manually}{%
\subsubsection{Run Task Manually}}

\emph{(Mã nguồn chi tiết đã được lược bỏ để báo cáo ngắn gọn. Vui lòng
xem trong source code dự án)}

\hypertarget{check-task-status}{%
\subsubsection{Check Task Status}}

\begin{Shaded}
\begin{Highlighting}[]
\CommentTok{\# List recent tasks}
\ExtensionTok{aws}\NormalTok{ ecs list{-}tasks }\AttributeTok{{-}{-}cluster} \VariableTok{$CLUSTER} \AttributeTok{{-}{-}family} \VariableTok{$TASK\_DEF} \AttributeTok{{-}{-}desired{-}status}\NormalTok{ STOPPED}

\CommentTok{\# Get task details}
\VariableTok{TASK\_ARN}\OperatorTok{=}\VariableTok{$(}\ExtensionTok{aws}\NormalTok{ ecs list{-}tasks }\AttributeTok{{-}{-}cluster} \VariableTok{$CLUSTER} \AttributeTok{{-}{-}family} \VariableTok{$TASK\_DEF} \AttributeTok{{-}{-}desired{-}status}\NormalTok{ STOPPED }\AttributeTok{{-}{-}query} \StringTok{"taskArns[0]"} \AttributeTok{{-}{-}output}\NormalTok{ text}\VariableTok{)}
\ExtensionTok{aws}\NormalTok{ ecs describe{-}tasks }\AttributeTok{{-}{-}cluster} \VariableTok{$CLUSTER} \AttributeTok{{-}{-}tasks} \VariableTok{$TASK\_ARN}

\CommentTok{\# View logs}
\ExtensionTok{aws}\NormalTok{ logs tail /ecs/newsrag{-}project }\AttributeTok{{-}{-}follow} \AttributeTok{{-}{-}filter{-}pattern} \StringTok{"crawler"}
\end{Highlighting}
\end{Shaded}

\hypertarget{monitoring-debugging}{%
\subsection{Monitoring \& Debugging}}

\hypertarget{cloudwatch-logs}{%
\subsubsection{CloudWatch Logs}}

\begin{Shaded}
\begin{Highlighting}[]
\CommentTok{\# Tail crawler logs}
\ExtensionTok{aws}\NormalTok{ logs tail /ecs/newsrag{-}project }\AttributeTok{{-}{-}follow} \AttributeTok{{-}{-}filter{-}pattern} \StringTok{"crawler"}

\CommentTok{\# Search for errors}
\ExtensionTok{aws}\NormalTok{ logs filter{-}log{-}events }\DataTypeTok{\textbackslash{}}
  \AttributeTok{{-}{-}log{-}group{-}name}\NormalTok{ /ecs/newsrag{-}project }\DataTypeTok{\textbackslash{}}
  \AttributeTok{{-}{-}filter{-}pattern} \StringTok{"ERROR"} \DataTypeTok{\textbackslash{}}
  \AttributeTok{{-}{-}start{-}time} \VariableTok{$(}\FunctionTok{date} \AttributeTok{{-}d} \StringTok{\textquotesingle{}1 hour ago\textquotesingle{}}\NormalTok{ +\%s}\VariableTok{)}\NormalTok{000}
\end{Highlighting}
\end{Shaded}

\hypertarget{cloudwatch-metrics}{%
\subsubsection{CloudWatch Metrics}}

Create a dashboard with: - \texttt{ECS/TaskCount} (Running, Pending,
Stopped) - \texttt{ECS/CPUUtilization} (per task) -
\texttt{ECS/MemoryUtilization} (per task) -
\texttt{SQS/NumberOfMessagesSent} (news\_raw queue) -
\texttt{SQS/NumberOfMessagesReceived} (Lambda consumer)

\hypertarget{common-issues}{%
\subsubsection{Common Issues}}

\begingroup
\small
\setlength{\tabcolsep}{3pt}
\renewcommand{\arraystretch}{1.15}
\begin{longtable}[]{@{}
  >{\raggedright\arraybackslash}p{(\columnwidth - 2\tabcolsep) * \real{0.4118}}
  >{\raggedright\arraybackslash}p{(\columnwidth - 2\tabcolsep) * \real{0.5882}}@{}}
\toprule\noalign{}
\begin{minipage}[b]{\linewidth}\raggedright
Issue
\end{minipage} & \begin{minipage}[b]{\linewidth}\raggedright
Solution
\end{minipage} \\
\midrule\noalign{}
\endhead
\bottomrule\noalign{}
\endlastfoot
\texttt{CannotPullContainerError} & Check ECR permissions, image tag
exists, VPC has Internet Gateway \\
\texttt{Task\ failed\ to\ start} & Check CloudWatch Logs for container
error; verify \texttt{command} in task definition \\
\texttt{SQS\ AccessDenied} & Attach SQS policy to
\texttt{ecs\_task\_role} \\
\texttt{Database\ connection\ timeout} & Verify security group allows
5432 from ECS SG; check Secrets Manager access \\
\texttt{Crawler\ takes\ \textgreater{}\ 15\ min} & Increase task
CPU/memory; optimize Scrapy settings (\texttt{CLOSESPIDER\_TIMEOUT}) \\
\end{longtable}
\endgroup

\hypertarget{scaling-considerations}{%
\subsection{Scaling Considerations}}

\begin{itemize}
\tightlist
\item
  \textbf{Concurrency}: Run multiple crawler tasks in parallel
  (different spiders)
\item
  \textbf{Queue depth}: Monitor SQS
  \texttt{ApproximateNumberOfMessagesVisible}
\item
  \textbf{Cost}: Fargate Spot (up to 70\% savings) for fault-tolerant
  crawling:
\end{itemize}

\begin{Shaded}
\begin{Highlighting}[]
\NormalTok{\# In task definition or capacity provider strategy}
\NormalTok{capacity\_provider\_strategy \{}
\NormalTok{  capacity\_provider = "FARGATE\_SPOT"}
\NormalTok{  weight            = 1}
\NormalTok{  base              = 0}
\NormalTok{\}}
\end{Highlighting}
\end{Shaded}

\hypertarget{updating-the-crawler}{%
\subsection{Updating the Crawler}}

\emph{(Mã nguồn chi tiết đã được lược bỏ để báo cáo ngắn gọn. Vui lòng
xem trong source code dự án)}

\hypertarget{next-steps}{%
\subsection{Next Steps}}

After crawler pushes to SQS: 1. \textbf{Lambda Consumer} reads SQS
\(\rightarrow\) inserts to Aurora: \href{5.11-Lambda-Consumer/}{Lambda
Consumer} 2. \textbf{Lambda ETL} processes raw \(\rightarrow\) Star
Schema + Bedrock Embed: \href{5.12-Lambda-ETL/}{Lambda ETL} 3.
\textbf{RAG API} serves queries: \href{5.13-RAG-API/}{RAG API}

\begin{center}\rule{0.5\linewidth}{0.5pt}\end{center}

\textbf{Next:} \href{5.11-Lambda-Consumer/}{Lambda Consumer (SQS
\(\rightarrow\) Aurora)}
